\chapter{Setting, conventions, and the three scores}
\label{ch:probability}
\label{ch:diffusion}

This chapter fixes what the rest of the document assumes. It is deliberately compact: the
material here is developed from first principles in the thesis (Hamiltonian and statistical
mechanics, stochastic calculus, the diffusion literature), and repeating that development would
turn a technical supplement into a second textbook. What is kept is what later chapters actually
cite --- the estimation-theoretic facts that Proposition~\ref{prop:closure-lmmse} rests on, the
coordinatewise structure of the channel that makes the whole project possible, and the
distinction between three different objects that the word ``score'' is used for.

\section{Conditioning and the Bayes-optimal estimator}
\label{sec:bayesopt}

Bayes' rule, in the form used throughout:
\begin{equation}
\label{eq:bayes}
P(a \mid x) \;=\; \frac{P(x\mid a)\,P(a)}{P(x)} \;\propto\;
\underbrace{P(x\mid a)}_{\text{likelihood}}\;\underbrace{P(a)}_{\text{prior}} .
\end{equation}

\begin{proposition}[Orthogonality of the posterior mean]
\label{prop:orthogonality}
For any estimator $u$ that is a function of $x$ alone,
\begin{equation}
\label{eq:orth}
\E\!\left[\|a-u\|^2 \,\middle|\, x\right]
\;=\;
\E\!\left[\|a-\pmean{a}\|^2 \,\middle|\, x\right]
\;+\;
\|u-\pmean{a}\|^2 .
\end{equation}
\end{proposition}

\begin{proof}
Expand $a-u = (a-\pmean{a}) + (\pmean{a}-u)$. The cross term vanishes because $\pmean{a}-u$ is
$x$-measurable and $\E[a-\pmean{a}\mid x]=0$.
\end{proof}

Two consequences are used constantly. The posterior mean minimises the conditional squared risk,
so it is the reference every estimator in this project is scored against. And the first term is
an irreducible floor common to all estimators, which is why every table reports error
\emph{against} $\pmean{a}$ rather than raw risk --- subtracting the floor is what makes the
comparison informative.

\begin{intuition}
Optimality here is with respect to squared error specifically. A different loss selects a
different functional of the posterior --- absolute error would select the median. Nothing in the
project depends on squared loss being the right choice for any application; it is the choice that
makes the target computable and the comparison well posed.
\end{intuition}

\section{Linear estimators}
\label{sec:lmmse-intro}

Restricting to estimators linear in $x$ gives the LMMSE estimator, which for zero-mean
$(a,x)$ is
\begin{equation}
\label{eq:lmmse-general}
\lmmse \;=\; \Cov(a,x)\,\Cov(x)^{-1} x .
\end{equation}

\begin{proposition}[Linear optimality]
\label{prop:lmmse-general}
Among estimators of the form $u = Bx$, \eqref{eq:lmmse-general} minimises $\E\|a-u\|^2$. When
$(a,x)$ is jointly Gaussian and zero-mean it coincides with $\pmean{a}$, since the conditional
expectation is then linear.
\end{proposition}

The zero-mean hypothesis is load-bearing rather than cosmetic: the general form carries
$\E[a] + \Cov(a,x)\Cov(x)^{-1}(x-\E[x])$, and dropping the mean terms when $\E[a]\neq0$ produces
a genuinely different estimator. Chapter~\ref{ch:gaussian} measures the size of that difference.

\section{Two things a distribution can be}
\label{sec:empirical}

A model is fitted to samples, not to a distribution. Writing the empirical measure
\begin{equation}
\label{eq:empirical}
\widehat{P}_0 \;=\; \frac{1}{\nseq}\sum_{\mu=1}^{\nseq} \delta_{a^\mu},
\end{equation}
the distinction between $\Pz$ and $\widehat{P}_0$ is the one \S\ref{sec:memorisation} turns into
the project's motivation. It is not a technicality about estimation error: the two have
genuinely different scores, and the reverse process driven by one does something categorically
different from the reverse process driven by the other.

\section{The forward process}

The Ornstein--Uhlenbeck channel,
\begin{equation}
\label{eq:ou}
\frac{\mathrm{d}x}{\mathrm{d}t} \;=\; -x + \sqrt{2}\,\gwn(t),
\qquad x(0) = a ,
\end{equation}
with $\gwn$ standardised Gaussian white noise, so the $\sqrt2$ is carried explicitly. Solving by
integrating factor gives
\begin{equation}
\label{eq:forward-solution}
x(t) \;=\; e^{-t} a + \sqrt{\Dt}\,z,
\qquad \Dt = 1 - e^{-2t},
\qquad z\sim\normal{\vect0}{\Id},
\end{equation}
and hence the noised distribution as a Gaussian convolution,
\begin{equation}
\label{eq:pt-def}
\Pt(x) \;=\; \int \mathrm{d}a\; \Pz(a)\,
\frac{1}{(2\pi\Dt)^{\nsites/2}}
\exp\!\left[-\frac{\|x - e^{-t}a\|^2}{2\Dt}\right].
\end{equation}
Two properties are used repeatedly: $\Dt\to0$ as $t\to0$, so $\Pt\to\Pz$; and $\Var(x_t)=1$ for
unit-variance data at every $t$, so initialising the reverse process at $\normal{\vect0}{\Id}$
incurs a bias of only $e^{-2\tmax}$.

\section{Why the coordinatewise structure matters}
\label{sec:coordinatewise}

This section is short and is the reason the project exists. The forward noise is independent
across coordinates, so the likelihood factorises sitewise:
\begin{equation}
\label{eq:likelihood-factorises}
P(x \mid a, t)
= \prod_{i=1}^{\nsites}
\frac{1}{\sqrt{2\pi\Dt}}
\exp\!\left[-\frac{(x_i - e^{-t}a_i)^2}{2\Dt}\right]
\;=\;\prod_{i=1}^{\nsites} \like(x_i \mid a_i).
\end{equation}

\begin{pitfall}
Each factor $\like(x_i\mid a_i)$ has degree one among the latent variables. Attaching a
degree-one factor to a graph adds a leaf and cannot create a cycle --- so conditioning on $x$
reweights node potentials without creating edges among the $a_i$, and the posterior factor graph
is whatever the prior's was.

Conditioning usually creates dependence. It does not here, precisely because each observation is
generated from a single latent variable. The construction fails the moment the forward noise is
correlated across coordinates: then $P(x\mid a)$ does not factorise, the likelihood becomes a
factor touching several $a_i$, and the posterior graph acquires edges the prior did not have.
Everything in Chapters~\ref{ch:bp}--\ref{ch:estimation} rests on this one property.
\end{pitfall}

\section{Reversing it}
\label{sec:reverse}

By the time-reversal theorem for diffusions \citep{anderson1982reverse,haussmann1986}, a forward SDE
$\mathrm{d}x = f\,\mathrm{d}t + g\,\mathrm{d}W$ has a time-reversed process with drift
$f - g^2\nabla\log p_t$ and the same diffusion coefficient. With $f=-x$ and $g=\sqrt2$, writing
$\tau = \tmax - t$ so that increasing $\tau$ means decreasing $t$,
\begin{equation}
\label{eq:reverse}
\frac{\mathrm{d}x}{\mathrm{d}\tau} \;=\; x + 2\,\score(x,\tau) + \sqrt{2}\,\gwn(\tau),
\qquad \score(x,t) := \nabla_x \log \Pt(x),
\end{equation}
which is the only object the reverse process needs. Chapter~\ref{ch:score} shows it is a
rescaling of the posterior mean, and therefore that generating and inferring are the same
problem.

\section{Map of the code}
\label{sec:codemap}

Everything referenced in this document lives under \codefile{research/nongaussian-bp/}.

\begin{center}
\small
\begin{tabular}{@{}p{5.3cm}p{9.3cm}@{}}
\toprule
\textbf{Path} & \textbf{Contents} \\
\midrule
\codefile{src/priors.py} & The data-generating models: \texttt{GaussianAR1},
\texttt{LaplaceAR1}, \texttt{StudentTAR1}, \texttt{UniformAR1},
\texttt{GaussianMixtureAR1}. Each provides \texttt{sample} and
\texttt{log\_transition\_matrix}. \\
\codefile{src/bp\_grid.py} & The inference recursion. \texttt{make\_grid},
\texttt{grid\_bp} (one sequence, with evidence), \texttt{grid\_bp\_batch} (many sequences,
device-parameterised). \\
\codefile{src/bp\_gaussian.py} & The closed-form Gaussian recursion in information form. \\
\codefile{src/exact\_scores.py} & Closed-form Gaussian posterior mean and score, used as
ground truth. \\
\codefile{src/em.py} & The expectation step, the sufficient statistic, and the estimation
driver. \\
\codefile{src/kernels.py} & Parametrised transition families and their maximisation steps. \\
\codefile{src/denoiser.py} & A common interface over inference-based and network-based
denoisers, plus the network training loop. \\
\codefile{src/reverse.py} & Reverse-time integration: SDE and probability-flow ODE. \\
\codefile{src/sample\_metrics.py} & Statistics for comparing generated data against the
truth. \\
\codefile{src/backend.py} & CPU/GPU array-module dispatch for the recursion. \\
\codefile{src/ring.py} & The rotating-ring model of Chapter~\ref{ch:ring}: gauge, the
closed-form $\psi$-conditional density, and the blind per-frame likelihood. \\
\codefile{experiments/frozen\_config.py} & The single configuration behind every number in the
note and the workshop paper. Imported, never re-specified. \\
\codefile{experiments/} & One file per experiment, each self-describing. \texttt{exp\_28} is
the ring. \\
\codefile{tools/check\_paper.sh} & The production gate: page limits, no code references, no
unfilled placeholders, shared sections still shared. \\
\codefile{tools/check\_replicates.py} & Parses each experiment's settings and fails on a
replicate count the frozen config cannot reach. See \S\ref{sec:corr-replicates}. \\
\codefile{tests/} & Validation. Several tests here are the evidence for claims made in this
document and are cited as such. \\
\bottomrule
\end{tabular}
\end{center}
